\begin{abstract}
Phase-coded frequency-modulated continuous-wave (PC-FMCW) signaling offers a natural way to combine radar sensing and communication, but high-mobility operation raises a question that is not addressed by fixed-waveform feasibility alone: which configuration should be used when the operating state is only imperfectly known? This paper studies that question for a \SI{77}{GHz} vehicular PC-FMCW physical layer. The proposed controller separates two decisions that are often conflated. It first removes waveform configurations that cannot support the requested range or radial velocity because of sampling and ambiguity limits. It then evaluates the remaining configurations under joint communication and sensing QoS constraints while accounting for uncertainty in SNR, synchronization residuals, interference, and related PHY state. Reliability is qualified with a finite-sample one-sided Wilson rule, and the controller abstains when no action is supported at the declared confidence level. In a frozen paired benchmark comprising 12,000 operating-state/seed pairs per policy, the robust policy selects 12.23\% of cases and succeeds in all 1,468 selected cases, giving a one-sided 95\% Wilson lower bound of 99.816\%. The deterministic joint policy selects more often (19.07\%) but achieves 84.83\% conditional joint QoS. Consequently, the robust policy has lower unconditional joint-QoS probability, and the paired B4--B3 difference is $-0.03942$ with a 95\% bootstrap interval of $[-0.04292,-0.03600]$. The main result is therefore not policy dominance, but a reliability--availability--computation trade-off: robust adaptation identifies a smaller operating region in which joint service is supported with stronger evidence, while explicitly abstaining elsewhere.
\end{abstract}
